% Generated draft: requires mathematical and source review.
\documentclass[11pt]{article}
\usepackage[margin=25mm]{geometry}
\usepackage{fontspec}
\setmainfont{Latin Modern Roman}
\usepackage{amsmath,amssymb,mathtools,unicode-math}
\setmathfont{Latin Modern Math}
\usepackage{enumitem,needspace,xurl,hyperref}
\setlength{\parindent}{0pt}
\setlength{\parskip}{5pt}
\setlength{\emergencystretch}{3em}
\newcommand{\N}{\mathbb{N}}
\newcommand{\Z}{\mathbb{Z}}
\newcommand{\Q}{\mathbb{Q}}
\newcommand{\R}{\mathbb{R}}
\newcommand{\C}{\mathbb{C}}
\newcommand{\F}{\mathbb{F}}
\newcommand{\A}{\mathbb{A}}
\newcommand{\PP}{\mathbb{P}}
\newcommand{\E}{\mathbb{E}}
\newcommand{\eps}{\varepsilon}
\newcommand{\dd}{\,\mathrm{d}}
\begin{document}
\section*{UnsolvedMath ID: 30003805}
\textbf{Source problem code:} catalogue entry
\subsection{Definitions and mathematical statement}
Under the usual complex algebraic convention, let \(A\) be a unital \(*\)-bialgebra over \(\C\), with coproduct \(\Delta:A\to A\otimes A\) and counit \(\varepsilon:A\to\C\). Here the tensor product is algebraic, the coproduct is coassociative, and \(\Delta\) and \(\varepsilon\) are unital \(*\)-homomorphisms satisfying the counit identities. No topology or analytic completion is specified in the supplied record. Put \(I=\ker\varepsilon\).

A generating functional is a complex-linear map \(L:A\to\C\) satisfying
\[
L(1)=0,\qquad L(a^*)=\overline{L(a)}\quad(a\in A),\qquad L(x^*x)\geq0\quad(x\in I).
\]
Write \(\mathcal G(A)\) for the set of these functionals. In the usual algebraic terminology, a Gaussian generating functional \(G\) satisfies
\[
G(abc)=0\qquad(a,b,c\in I).
\]
These definitions use the standard complex algebraic convention; the supplied record does not explicitly specify the scalar field or an analytic setting.

\textbf{Requested classification, with an unresolved definition.} Determine which unital \(*\)-bialgebras have the L\'evy--Khintchine property: every generating functional admits a L\'evy--Khintchine-type decomposition. The supplied record does not define the admissible components or conditions of this decomposition. To display the missing information explicitly, let \(\mathcal D(A)\subseteq\mathcal G(A)\times\mathcal G(A)\) denote the pairs admissible under the intended decomposition convention. The requested property would then be
\[
\forall L\in\mathcal G(A)\quad\exists (G,J)\in\mathcal D(A)\quad\forall a\in A,\qquad L(a)=G(a)+J(a).
\]
A Gaussian-plus-jump formulation would require \(G\) to be Gaussian and would need a precise condition on \(J\), together with any required compatibility or maximality condition. Those conditions are not supplied. Merely allowing arbitrary generating functionals as the second component would make existence trivial by taking \(G=0\) and \(J=L\). Thus \(\mathcal D(A)\) remains an explicit unresolved part of the intended statement, and the display is a specification of the classification target rather than a complete definition of the L\'evy--Khintchine property.

The original wording additionally singles out \(SU_q(d)\) for integers \(d\geq3\) as unresolved examples. It specifies neither the allowed values of \(q\) nor the associated function or coordinate \(*\)-bialgebra. The supplied later literature assessment instead reports the property for all \(SU_q\) and \(U_q\) quantum groups, without parameter or algebra conventions. Preserve the former as a historical target and the latter as supplied dated curation; their precise scope and relationship require review.
\subsection{Short English statement}
Classify the unital \(*\)-bialgebras for which every generating functional has the intended L\'evy--Khintchine decomposition. The supplied record leaves that decomposition undefined, so its admissible components must be clarified before the classification has a fully precise target. The original \(SU_q(d)\), \(d\geq3\), examples also lack parameter and algebra conventions. The supplied partially solved label and later claims about quantum groups are dated curation, not independently verified conclusions.
\subsection{Sources}
\begin{itemize}
\item Local UnsolvedMath registry, supplied record with permanent ID \(30003805\), title \emph{The L\'evy--Khintchine Property for Quantum Groups}, source problem code \texttt{catalogue entry}, in \emph{Oberwolfach Reports Open Problems}. The record gives proposed year 2018 and names no individual proposer.
\item Supplied report attribution: \emph{Interactions between Operator Space Theory and Quantum Probability with Applications to Quantum Information} (2019), Oberwolfach Reports, pp.\ 1293--1379, DOI \texttt{10.4171/owr/2018/22}. \url{https://doi.org/10.4171/owr/2018/22}. The background describes this report as documenting workshop open-problem discussions. No exact page or passage identifying this problem is supplied.
\item Supplied literature source, labelled \emph{Hunt formula for SUq and Uq}: \url{https://eprints.lancs.ac.uk/id/eprint/222139}.
\item Supplied literature source, labelled \emph{Generating functional decompositions}: \url{https://arxiv.org/abs/1510.03292}.
\item Supplied literature source, labelled \emph{Oberwolfach Report 22}: \url{https://doi.org/10.4171/owr/2018/22}.
\item Literature assessment dated exactly \texttt{2026-08-21}: the supplied status is \emph{partially solved}; the supplied assessment reports the LK property for broad classes and all \(SU_q\) and \(U_q\) quantum groups, while reporting counterexamples and an outstanding general classification. This is supplied dated curation derived from the OWR record, not independently verified here. Its time-sensitive status and cited claims require independent verification.
\end{itemize}
\end{document}
